\section{Birational decompositions and the geometric reduction}
\label{sec:inductive-reduction}

The proof will produce a semiample line bundle on a smooth model, together
with its entire fixed divisorial part. This section makes that statement
precise and reduces the induction to two cases: spaces admitting a
nontrivial fibration, and simple spaces of algebraic dimension zero.

\subsection{The inductive decomposition and its negative part}

We use additive notation for rational holomorphic line bundles. Thus
\(L\sim_{\Q}L'\) means that \(mL\) and \(mL'\) are holomorphically
isomorphic for some positive integer \(m\). If a rational divisor occurs in such
an identity, it denotes its associated rational line bundle. In contrast,
\(c_1(L)\) and \(\{D\}=c_1(\OO_X(D))\) denote real Bott--Chern classes.
Canonical comparisons are made with the usual local meromorphic canonical
identifications. In particular, an identity of lines below contains more
information than equality of their Chern classes.

Let \(X\) be a smooth compact Kähler manifold and let \(\alpha\) be a
pseudo-effective real \((1,1)\)-class. Fix a Kähler form \(\omega\). For a
prime divisor \(D\), its minimal multiplicity is
\[
 \nu_D(\alpha)=\lim_{\varepsilon\downarrow0}
 \inf\bigl\{\nu_D(T):T\in\alpha,\ T\geq-\varepsilon\omega\bigr\}.
\]
Here \(T\) ranges over closed currents, and \(\nu_D(T)\) is its generic
Lelong number; one may equivalently use currents with analytic singularities.
The limit is independent of \(\omega\). Boucksom's divisorial decomposition
is
\[
 N(\alpha)=\sum_D\nu_D(\alpha)D,
 \qquad Z(\alpha)=\alpha-\{N(\alpha)\}.
\]
The sum is a finite effective real divisor, and \(Z(\alpha)\) is modified
nef: its minimal multiplicity at every prime is zero. We write \(N(L)\) and
\(\nu_D(L)\) for \(N(c_1(L))\) and \(\nu_D(c_1(L))\). These are the
analytic notions, with small Kähler
perturbations, throughout the paper. We use the foundational results in
\cite[Sections~2--3 and~5]{Boucksom2004}.

The induction asks for a decomposition that retains both the actual line
bundle and this analytic negative divisor.

\begin{definition}\label{def:good-decomposition}
For \(n\geq0\), let \(\Ggood_n\) be the following assertion. If \(X\)
is a connected smooth compact Kähler manifold of dimension \(n\), \(B\)
is a rational simple normal crossing boundary with coefficients in
\([0,1]\), and \(J=K_X+B\) is pseudo-effective, then there is a smooth
compact Kähler modification \(\mu:Y\to X\) and an actual rational line
identity
\begin{equation}\label{bir:good-decomposition}
 \mu^*J\sim_{\Q}P+R,
 \qquad P\text{ semiample},\qquad R=N(\mu^*J)\geq0,
\end{equation}
where \(R\) is a rational divisor.
\end{definition}

The next lemmas make this decomposition stable under further resolutions
and allow exceptional resolution errors to be removed. They also provide
the fixed-section statement needed to descend the case in which the
positive part is torsion.

\begin{lemma}[Negative-part calculus]\label{bir:calculus}
Let \(\alpha\) be pseudo-effective on a smooth compact Kähler manifold.
\begin{enumerate}
\item Every positive current in \(\alpha\) contains the divisorial current
\([N(\alpha)]\). If \(0\leq F\leq N(\alpha)\), then
\[
 \alpha-\{F\}\text{ is pseudo-effective},\qquad
 N(\alpha-\{F\})=N(\alpha)-F.
\]
\item Minimal multiplicities are homogeneous and subadditive. In
particular, if \(\beta\) is nef, then
\(N(\alpha+\beta)\leq N(\alpha)\).
\item If \(\mu:Y\to X\) is a smooth modification and \(D'\) is the strict
transform of a prime \(D\), then
\[
 \nu_{D'}(\mu^*\alpha)=\nu_D(\alpha).
\]
\item If \(\gamma\) is modified nef and \(j:\widetilde D\to X\) is a
resolution of a prime divisor, then \(j^*\gamma\) is pseudo-effective.
\end{enumerate}
\end{lemma}

\begin{proof}
The first assertion about currents follows from the definition of minimal
multiplicity and Siu decomposition. Homogeneity and subadditivity follow by
scaling and adding testing currents. The subtraction identity is the
corresponding property in the big cone, followed by a small Kähler
perturbation. More explicitly, for \(\alpha_\varepsilon=\alpha+
\varepsilon\{\omega\}\), subtract
\(F_\varepsilon=\min\{F,N(\alpha_\varepsilon)\}\) coefficientwise.
In the big cone all positive currents contain this divisor, so subtracting
it translates each minimal multiplicity by its coefficient. As
\(\varepsilon\downarrow0\), \(F_\varepsilon\to F\). The class of
\(F-F_\varepsilon\) can be absorbed in a Kähler perturbation tending to
zero (there are only finitely many components). This gives the upper bound
for the asserted identity; subadditivity applied after adding \(F\) gives
the reverse bound. Weak compactness of positive currents gives
pseudo-effectivity of the limit. This is also the subtraction property of
the divisorial decomposition in \cite[Section~3]{Boucksom2004}.

For the strict-transform formula, pull almost-positive testing currents
to \(Y\). Near the generic point of \(D'\), the map is an isomorphism,
so the generic order is unchanged. This proves one inequality. Conversely,
push a testing current on \(Y\) to \(X\). To control its negative error,
write the error as a small multiple of the fixed positive current
\(\mu_*\omega_Y\), and add a fixed smooth Kähler representative of
\(C\omega_X-\{\mu_*\omega_Y\}\) for \(C\) sufficiently large. The
result is a positive test in a Kähler perturbation tending to zero. The
pushforward of \(\omega_Y\) has no divisorial order at the generic point
of \(D\), where \(\mu\) is an isomorphism. Thus this test has the same
generic order as the original one at \(D'\). Regularization, if needed,
returns to tests with analytic singularities. This proves the other
inequality and the formula. It is the analytic form of the strict-transform
comparison in \cite[Lemma~5.8]{Boucksom2004}.

Finally, choose analytic-singularity tests for a modified-nef class whose
generic order at the chosen prime tends to zero. Subtract that generic
divisorial order before restricting to a resolution of the prime. The
remaining singular set does not contain its generic point, so restriction
gives an almost-positive current there. The subtracted multiple and the
Kähler error tend to zero. Taking the limit proves the last assertion; see
also \cite[Propositions~2.4 and~3.8]{Boucksom2004}.
\end{proof}

For a rational line \(J\) on a normal compact Kähler space \(X\) whose pullback
to a smooth resolution is pseudo-effective, we also use minimal
multiplicities for divisorial places. If \(Q\) appears on a smooth
resolution \(p:Y\to X\), set
\[
 \nu_Q(J):=\nu_Q(p^*J).
\]
The strict-transform formula on a common smooth refinement makes this
independent of the chosen resolution. When a real class is already on a
fixed smooth model, \(\nu_Q\) continues to denote its coefficient in the
negative part on that model.

\begin{lemma}[Exceptional translation]\label{bir:exceptional}
Let \(p:Y\to X\) be a proper bimeromorphic morphism from a smooth compact
Kähler manifold to a normal compact Kähler space. Let \(M\) be a real
\((1,1)\)-class on \(X\) represented by smooth local potentials, and let
\(E\geq0\) be a real \(p\)-exceptional divisor. If
\(\alpha=p^*M+\{E\}\) is pseudo-effective, then \(p^*M\) is
pseudo-effective and
\[
 N(\alpha)=N(p^*M)+E.
\]
\end{lemma}

\begin{proof}
It suffices by Lemma~\ref{bir:calculus} to prove \(E\leq N(\alpha)\).
Write
\[
 E-N(\alpha)=U-V,
 \qquad U,V\geq0,
\]
with no common component, and suppose \(U\neq0\). Its components are
\(p\)-exceptional. Put \(n=\dim Y\) and
\(\ell=\max\{\dim p(U_i):U_i\text{ a component of }U\}\). Normality
gives \(\ell\leq n-2\). Let \(\eta\) be the pullback of a Kähler form
on \(X\), let \(\omega\) be Kähler on \(Y\), and set
\[
 q(\gamma,\delta)=
 \int_Y\gamma\,\delta\,\eta^{\ell}\omega^{n-\ell-2}.
\]
The restriction property in Lemma~\ref{bir:calculus} gives
\(q(Z(\alpha),U)\geq0\). The term \(q(p^*M,U)\) is zero: on each
component of \(U\), it contains \(\ell+1\) factors pulled back from an
image of dimension at most \(\ell\). Distinct effective divisors have
nonnegative intersection against the semipositive and Kähler factors in
\(q\). Hence
\[
 0\leq q(Z(\alpha),U)=q(U,U)-q(V,U)\leq q(U,U).
\]
On the other hand, \(q(U,\eta)=0\) by the same dimension count, whereas
\(q(\eta,\eta)>0\). The mixed Hodge index theorem, applied first with
Kähler factors and then by a semipositive limit, says that \(q\) is
negative semidefinite on \(\eta^\perp\). It follows that \(q(U,U)=0\)
and that \(U\) is in the radical of \(q\): the radical assertion follows
from negative semidefiniteness on \(\eta^\perp\), and then from
\(q(U,\eta)=0\). But
\[
 q(U,\omega)=\sum_i u_i\int_{U_i}
       \eta^{\ell}\omega^{n-\ell-1}>0.
\]
Indeed, a component with image dimension \(\ell\) has strictly positive
integrand on a dense open. This is a contradiction. The mixed Hodge index
statement used here is the mixed Hodge--Riemann relation
\cite[Theorems~A and~C]{DinhNguyen}; the limit preserves the assertion that
there is at most one positive direction.
\end{proof}

\begin{lemma}[Pulling up a known decomposition]\label{bir:pullup}
Let \(\alpha\) be a pseudo-effective real \((1,1)\)-class on a smooth compact
Kähler manifold, and suppose
\[
 \alpha=\beta+\{R\},\qquad \beta\text{ nef},\qquad R=N(\alpha)\geq0.
\]
For every smooth modification \(\mu:Y\to X\),
\[
 N(\mu^*\alpha)=\mu^*R.
\]
In particular, this applies to the Chern classes of an actual rational line
decomposition with a nef rational positive part.
\end{lemma}

\begin{proof}
Subadditivity and nefness give \(N(\mu^*\alpha)\leq\mu^*R\). By the
strict-transform formula, the difference
\(U=\mu^*R-N(\mu^*\alpha)\) is effective and exceptional. Moreover
\(Z(\mu^*\alpha)=\mu^*\beta+\{U\}\) is modified nef. Apply
Lemma~\ref{bir:exceptional} to this class. It gives
\(U\leq N(Z(\mu^*\alpha))=0\), proving the assertion.
\end{proof}

\begin{lemma}[Fixed sections]\label{bir:sections}
Let \(L\) be a pseudo-effective rational line on a smooth compact Kähler
manifold. Every section \(s\) of an integral multiple \(mL\) satisfies
\[
 \ord_D(s)\geq m\nu_D(L)
\]
at every prime \(D\). If \(L\sim_{\Q}P+R\), with \(P\) semiample and
\(R=N(L)\) rational, then in every sufficiently divisible degree
multiplication by the canonical section of \(mR\) identifies
\(H^0(X,mP)\) with \(H^0(X,mL)\).
\end{lemma}

\begin{proof}
The divisor current \([\operatorname{div}(s)]/m\) is a positive current
in \(c_1(L)\), so Lemma~\ref{bir:calculus} gives the order bound. In the
stated decomposition, every section therefore divides holomorphically by
the canonical section of \(mR\). Conversely, multiplication gives a
section of \(mL\). The actual rational line identity makes these operations
inverse after clearing denominators.
\end{proof}

\begin{lemma}[Uniqueness when the positive part is torsion]
\label{bir:unique-current}
Suppose \(L\sim_{\Q}R\), where \(R=N(L)\) is an effective rational
divisor on a smooth compact Kähler manifold. Then the only positive current
in \(c_1(L)\) is \([R]\). If \(F\geq0\) is a rational divisor and
\(L-F\) is pseudo-effective, then
\[
 F\leq R,\qquad L-F\sim_{\Q}R-F=N(L-F).
\]
The same conclusions hold if \(L\sim_{\Q}P+R\) with \(P\) torsion.
\end{lemma}

\begin{proof}
Write \(X\) for the manifold. The assertion is immediate if \(\dim X=0\),
so assume \(\dim X>0\).
Every positive current \(T\) in \(c_1(L)\) contains \([R]\). The
residual \(T-[R]\) is positive with zero cohomology class, so its mass
against a Kähler form to the power \(\dim X-1\) is zero. Thus it vanishes.
For the second assertion, add \([F]\) to any positive current in
\(c_1(L-F)\). Uniqueness gives \(F\leq R\), and subtraction in
Lemma~\ref{bir:calculus} gives the asserted negative part. A torsion line
has zero Chern class and becomes trivial after taking a positive multiple,
so the last statement is the same argument.
\end{proof}

\subsection{Birational transfer of the decomposition}

All modifications resolving spaces, maps, or ideals can be chosen
projective. A projective modification of a compact Kähler space is Kähler;
graphs between compact Kähler models can likewise be resolved by smooth
compact Kähler manifolds. We use these standard analytic resolution and
flattening results without further mention.

For log discrepancies our convention is
\[
 a(E;X,\Delta)=1+\ord_E\bigl(K_Y-p^*(K_X+\Delta)\bigr)
\]
on a smooth model \(p:Y\to X\). In particular, on a log resolution of an
lc pair, giving each new exceptional divisor coefficient one produces an
effective exceptional error. The next proposition explains why this
convention is harmless.

\begin{proposition}[Resolution and descent]\label{bir:resolution-transfer}
Let \((X,\Delta)\) be a normal compact Kähler lc pair with rational
boundary and rational Cartier adjoint \(J\). Let \(p:Y\to X\) be a log
resolution, and put
\[
 B_Y=p_*^{-1}\Delta+\sum_{E\text{ exceptional}}E.
\]
If \(p^*J\) is pseudo-effective and the conclusion of
Definition~\ref{def:good-decomposition} holds for \((Y,B_Y)\), then there
is a smooth compact Kähler modification \(r:V\to X\) with an actual
rational line identity
\[
 r^*J\sim_{\Q}P+R,\qquad P\text{ semiample},\qquad
 R=N(r^*J)\geq0,
\]
where \(R\) is a rational divisor.
If \(J\) is analytically nef, it is semiample on \(X\).
\end{proposition}

\begin{proof}
With compatible canonical choices there is an actual rational divisor
identity
\[
 K_Y+B_Y=p^*J+E_p,\qquad
 E_p=\sum_{E\text{ exceptional}}a(E;X,\Delta)E\geq0.
\]
Suppose \(q:V\to Y\) realizes the decomposition for the left side.
Lemma~\ref{bir:exceptional}, applied over \(X\), gives
\[
 N\bigl(q^*(K_Y+B_Y)\bigr)
 =N(q^*p^*J)+q^*E_p.
\]
Cancelling \(q^*E_p\) in the actual rational line identities proves the
first assertion. This argument also shows that further resolutions using
the reduced-exceptional convention do not change the assertion
\(\Ggood_n\).

If \(J\) is nef, its pullback is nef, so its negative part is zero. The
semiample line in the decomposition is therefore \(q^*p^*J\) itself, up
to rational line isomorphism. Choose an actual Cartier multiple generated
on \(V\). Normality gives \((p q)_*\OO_V=\OO_X\), and projection formula
identifies its sections with the sections downstairs. A base point
downstairs would make every pullback section vanish on its nonempty fiber.
There is no such point, so that multiple of \(J\) is globally generated.
\end{proof}

The second transfer concerns a torsion positive part. Its proof uses the
following local construction for finite maps, which will also be used for
meromorphic sections and for covers of spaces of algebraic dimension zero.

\begin{lemma}[Finite analytic norms and characteristic polynomials]
\label{bir:finite-norm}
Let \(\nu:Z\to X\) be a finite surjective morphism of normal irreducible
complex spaces, of degree \(d>0\), and let \(L\) be a holomorphic line
bundle on \(X\). A nonzero meromorphic section \(s\) of \(\nu^*L\)
has a nonzero meromorphic norm \(\operatorname{Nm}_\nu(s)\) in
\(L^{\otimes d}\), with
\[
 \operatorname{div}(\operatorname{Nm}_\nu(s))
   =\nu_*\operatorname{div}(s).
\]
If \(s\) is holomorphic, its norm is holomorphic. If, in addition,
\(s\) is nonzero at every point of \(\nu^{-1}(x)\), then its norm is
nonzero at \(x\).

For a global meromorphic function \(f\) on \(Z\), there is a monic
polynomial \(\chi_f(T)\) of degree \(d\) with global meromorphic
coefficient functions on \(X\) such that \(\chi_f(f)=0\) after pulling
the coefficients to \(Z\).
\end{lemma}

\begin{proof}
At \(x\in X\), choose a holomorphic frame of \(L\) and put
\[
 R=\OO_{X,x},\qquad K=\operatorname{Frac}(R),\qquad
 A=(\nu_*\OO_Z)_x.
\]
The generic algebra \(A\otimes_RK\) is a product of finite field
extensions of \(K\), of total dimension \(d\). The local meromorphic
coefficient \(a\) of \(s\) belongs to this algebra. Since \(Z\) is
irreducible and \(s\) is not identically zero, its restriction is nonzero
on every generic local component. Thus \(a\) is nonzero in every field
factor, and the determinant of multiplication by \(a\) is nonzero.
Under a change of frame by a unit \(g\), this determinant changes by
\(g^{-d}\). The determinants therefore glue to the asserted
meromorphic section. The valuation formula for a norm over a discrete
valuation ring gives the displayed divisor identity at every prime of
the normal space \(X\).

If \(s\) is holomorphic, then \(a\in A\) is integral over \(R\).
Its norm is integral over \(R\) and belongs to \(K\), hence belongs to
\(R\) by normality. This does not require \(A\) to be locally free.
If \(s\) is nonzero at every point over \(x\), then \(a\) is a unit
in the finite semilocal algebra \(A\); applying the same argument to
\(a^{-1}\) shows that its norm is a unit in \(R\).

For a meromorphic function \(f\), use instead the determinant of
\(T\) minus multiplication by its local coefficient. These monic
polynomials agree on overlaps, and Cayley--Hamilton shows that they
annihilate \(f\). All these constructions use fractions of local analytic
germs, so they remain available when the global meromorphic functions on
\(X\) are constant.
\end{proof}

We now descend a decomposition whose positive part is torsion, the form
needed for a covering family on a space of algebraic dimension zero.

\begin{proposition}[Descent of a purely negative decomposition]
\label{bir:finite-descent}
Let \(e:Y\to X\) be a proper generically finite surjective morphism
between connected smooth compact Kähler manifolds. Let \(L\) be a
pseudo-effective rational line on \(X\), and let \(F\geq0\) be a
rational divisor on \(Y\). Suppose that on a smooth model over \(Y\),
the line \(e^*L+F\) is rationally linearly equivalent to its rational
negative part. Then there is an effective rational divisor \(D\) on
\(X\) such that
\[
 L\sim_{\Q}D=N(L).
\]
\end{proposition}

\begin{proof}
Replace \(Y\) by the smooth model in the hypothesis and pull \(F\) back.
Write \(e^*L+F\sim_{\Q}R=N(e^*L+F)\). Pull back a positive current in
\(c_1(L)\) and add \([F]\). Lemma~\ref{bir:unique-current} gives
\(F\leq R\). Put \(D_Y=R-F\). Subtraction in
Lemma~\ref{bir:calculus} gives
\[
 e^*L\sim_{\Q}D_Y=N(e^*L)\geq0.
\]
In particular \([D_Y]\) is the unique positive current in \(c_1(e^*L)\).

Choose \(m\) so that \(mL\) and \(mD_Y\) are integral and the latter is
the divisor of a holomorphic section \(s\) of \(e^*(mL)\). Factor \(e\)
through its normal Stein space:
\[
 Y\xrightarrow{h}Z\xrightarrow{\nu}X,
 \qquad \nu\text{ finite of degree }d=\deg e.
\]
The connected birational map \(h\) descends \(s\) to a holomorphic
section \(s_Z\) of \(\nu^*(mL)\), by normality and projection formula.
Lemma~\ref{bir:finite-norm} gives a holomorphic section of \(mdL\) whose
divisor is
\[
 \operatorname{div}(\operatorname{Nm}_\nu(s_Z))
   =\nu_*\operatorname{div}(s_Z)=m e_*D_Y.
\]
Thus \(D=d^{-1}e_*D_Y\) is effective and satisfies \(L\sim_{\Q}D\).

Both \(e^*D\) and \(D_Y\) are positive divisor currents in
\(c_1(e^*L)\), so uniqueness gives \(e^*D=D_Y\). We already know
\(N(L)\leq D\). If this inequality were strict at a prime \(A\subset X\),
choose analytic-singularity tests for \(L\) whose generic orders at \(A\)
tend to \(\nu_A(L)\), and pull them to \(Y\). For any prime \(A'\)
dominating \(A\), the generic orders are multiplied by
\(\operatorname{mult}_{A'}e^*A\). The pulled-back Kähler errors can be
enlarged to Kähler errors upstairs tending to zero. Consequently
\[
 \nu_{A'}(e^*L)
 \leq \operatorname{mult}_{A'}(e^*A)\,\nu_A(L)
 <\operatorname{mult}_{A'}(e^*A)\,\operatorname{coeff}_A(D),
\]
contrary to \(N(e^*L)=D_Y=e^*D\). This proves \(D=N(L)\).
\end{proof}

We will also use Proposition~\ref{bir:finite-descent} when \(e\) is a
modification and \(F=0\). Once the positive part of a decomposition is
known to be torsion, it gives \(L\sim_{\Q}N(L)\) on the original smooth
space; Lemma~\ref{bir:pullup} then identifies the negative divisor upstairs
with the pullback of \(N(L)\).

\subsection{The projective anchor and the two geometric cases}

\begin{proposition}[Projective good models]\label{proj:good-models}
Under Assumption~\ref{asm:log-iitaka}, every projective lc pair over
\(\C\) with rational boundary and pseudo-effective rational Cartier
adjoint has a good log minimal model. In particular, \(\Ggood_n\) holds
for projective manifolds in every dimension, and every nef rational lc
adjoint on a projective variety is semiample as an actual rational line.
\end{proposition}

\begin{proof}
The good-model assertion is the rational-boundary consequence of the full
argument in \cite[Proposition~2.5, Theorem~9.6, and Section~10]{LI}. That
argument establishes the additional inductive and nonvanishing premises of
its Proposition~2.5 before applying it, and proves a stronger real-boundary
induction. Its only use of Assumption~\ref{asm:log-iitaka} is in its
Lemma~6.1, on a resolved projective Albanese fibration with both boundaries
zero. Thus its premise is precisely available here.

For clarity, its good-model conclusion gives on a common smooth projective
resolution
\[
 u^*(K_X+B)=v^*(K_{X_m}+B_m)+F,
 \qquad F\geq0\text{ and }v\text{-exceptional},
\]
as an actual rational divisor identity; the line on \(X_m\) is semiample.
This is the comparison in \cite[Lemma~2.2]{LI}. Algebraic nefness of a
rational line on a projective variety implies analytic nefness, by adding
arbitrarily small ample rational classes. Lemma~\ref{bir:exceptional}
therefore identifies \(F\) with the analytic negative part of the left
side. On a smooth projective variety, analytic pseudo-effectivity of a
divisor class agrees with algebraic pseudo-effectivity
\cite[Proposition~1.2]{BDPP}; hence the good-model assertion applies to
every projective instance of \(\Ggood_n\). This proves the stated form of
\(\Ggood_n\). If the original adjoint
is nef, the same comparison and normal descent give its semiampleness.
\end{proof}

The remainder of the paper proves the following two propositions. They
are stated here so that the global induction can be completed before its
technical components are developed. A compact space is \emph{simple} if
no positive-dimensional proper compact analytic subvariety passes through
a very general point. We write \(a(X)\) for its algebraic dimension.

\begin{proposition}[The fibration case]\label{rank:fibration-case}
Assume Assumption~\ref{asm:log-iitaka}. Fix \(n>0\) and assume
\(\Ggood_j\) for \(j<n\). Let \(X\) be a connected
smooth compact Kähler manifold of dimension \(n\), let \(B\) be a rational
SNC boundary with coefficients in \([0,1]\), and suppose \(K_X+B\) is pseudo-effective. If \(X\) admits
a dominant meromorphic map to an irreducible compact space \(Y\) in Fujiki
class \(\mathcal C\), with \(0<\dim Y<n\), then the conclusion of
\(\Ggood_n\) holds for \((X,B)\).
\end{proposition}

\begin{proposition}[The simple case]\label{simple:case}
Fix \(n>0\) and assume \(\Ggood_j\) for \(j<n\). Let \(X\) be a connected
smooth simple compact Kähler manifold of dimension \(n\) with \(a(X)=0\).
For a rational SNC boundary \(B\) with coefficients in \([0,1]\) and
\(J=K_X+B\) pseudo-effective,
there is a smooth compact Kähler modification \(\mu:Y\to X\) such that
\[
 \mu^*J\sim_{\Q}N(\mu^*J),
\]
and this negative part is a rational divisor. In particular, the conclusion
of \(\Ggood_n\) holds with torsion positive part.
\end{proposition}

Proposition~\ref{rank:fibration-case} is proved in
Section~\ref{sec:rank-one}; it reduces the fibration to relative log Kodaira
dimension zero and then descends the adjoint to its base.
Proposition~\ref{simple:case} is proved in Section~\ref{sec:signed}, using
the meromorphic nonvanishing theorem of Section~\ref{sec:nonvanishing}.
Both propositions use \(\Ggood\) only in dimensions below \(n\).

\begin{theorem}[Good divisorial decomposition]
\label{thm:good-decomposition}
Under Assumption~\ref{asm:log-iitaka}, the assertion \(\Ggood_n\) holds
for every finite \(n\).
\end{theorem}

\begin{proof}
We use Propositions~\ref{rank:fibration-case} and~\ref{simple:case}, whose
proofs occupy the rest of the paper. The assertion is immediate in dimension
zero. Assume it in smaller dimensions and let \((X,B)\) be as in
Definition~\ref{def:good-decomposition} in dimension \(n>0\).

If \(a(X)=n\), the Kähler Moishezon theorem makes \(X\) projective, and
Proposition~\ref{proj:good-models} applies. If \(0<a(X)<n\), the algebraic
reduction of \(X\) is a dominant meromorphic map to a projective model of
dimension \(a(X)\) \cite[Theorem and Definition~3.1]{CampanaPeternell1999}.
Proposition~\ref{rank:fibration-case} applies.

It remains to consider \(a(X)=0\). If \(X\) is simple, use
Proposition~\ref{simple:case}. Otherwise Campana's maximal covering-family
theorem supplies a generically finite evaluation map from an incidence
space which itself has a nontrivial fibration; see
\cite[Lemma~17 and Corollary~18]{Campana2026}. After resolving the incidence,
the graph, and the base, we obtain a smooth compact Kähler source \(X'\) and maps
\[
 e:X'\longrightarrow X,\qquad f:X'\longrightarrow Y',
 \qquad 0<\dim Y'<n,
\]
where \(e\) is proper generically finite and \(Y'\) is in class
\(\mathcal C\) \cite[(1.1.1), (1.5)(3)--(4), (2.2)(1),(3), and (2.3)]{Fujiki1982}.
The algebraic dimension of \(X'\) is zero. Indeed, a
meromorphic function on \(X'\) descends through the birational part of the
normal Stein factorization of \(e\). Lemma~\ref{bir:finite-norm} gives its
characteristic polynomial over the finite part, with meromorphic
coefficients on \(X\). These coefficients are constant because \(a(X)=0\),
so the function is constant on the irreducible space \(X'\). Bimeromorphic
modifications do not change this conclusion.

Choose a reduced SNC divisor \(B'\) containing the inverse image of
\(\Supp B\), after a further resolution. Pullback of logarithmic
differentials gives the actual rational identity
\[
 K_{X'}+B'=e^*(K_X+B)+F,\qquad F\geq0.
\]
For example, this follows locally by pulling back logarithmic top forms
for the reduced support of \(B\); decreasing its coefficients to those of
\(B\) only increases the error. Thus the adjoint upstairs is
pseudo-effective. Proposition~\ref{rank:fibration-case} applies to
\((X',B')\). Its semiample positive part is torsion, since a nonconstant
semiample map would give a nonconstant meromorphic function on \(X'\).
Proposition~\ref{bir:finite-descent} now gives the required decomposition
for \(K_X+B\). This exhausts the cases and proves the induction.
\end{proof}

\begin{proof}[Proof of Theorem~\ref{thm:main}]
Apply Theorem~\ref{thm:good-decomposition} on a log resolution of the
given lc pair, with the reduced-exceptional boundary. The adjoint is
pseudo-effective because the original one is nef.
Proposition~\ref{bir:resolution-transfer} then gives generation of an
actual Cartier multiple at every point of the original normal space.
\end{proof}
